\documentclass[11pt]{article}
\usepackage{fontspec}
\usepackage[french]{babel}
\usepackage{amsmath}
\usepackage{amssymb}
\usepackage{amsthm}
\usepackage{mathrsfs}
\usepackage{tikz-cd}
\usepackage{enumitem}
\usepackage{lmodern}

\newcommand{\printedpage}[1]{\marginpar{\footnotesize p.~#1}}
\newcommand{\msp}[1]{\marginpar{\footnotesize [#1]}}
\newcommand{\hot}{\mathrm{hot}}
\newcommand{\HOT}{\mathrm{HOT}}
\newcommand{\Cat}{\mathrm{Cat}}
\newcommand{\Liss}{\underline{\mathrm{Liss}}}
\newcommand{\multhot}{\mathrm{multhot}}
\newcommand{\mult}{\mathrm{mult}}

\begin{document}

\printedpage{94}
% head: A. Grothendieck : Dérivateurs, Ch. VI, Édition des Universités de Montpellier II et Paris VII
Si p.ex. $s$ est élément final $e$ de $S$, on trouve
\[\Phi(\underset{S}{e}) \simeq S,\]
donc
\[\hot_{S,W}(X)\times_S \hot_{S,W}(e)=\hot_{S,W}(X),\]
d'autre part
\[\hot_{S,W}(X\times_S e)=\hot_{S,W}(X_e).\]

\msp{154}
Si les deux étaient isomorphes dans $\HOT_W(S)$ par la flèche canonique
\[X_e\longrightarrow X,\]
en appliquant le foncteur $p_!$ (où $p : S\longrightarrow e$), il s'ensuiverait que $X_e\longrightarrow X$ est dans $W$. Si par exemple $S=\mathbf{\Delta}^1$, cela impliquerait que pour tout $X$ lisse sur $S$, $X_1\longrightarrow X$ est dans $W$, alors que normalement la lissité implique que $X_0\longrightarrow X$ est dans $W$. Donc ça reviendrait à dire que $\Liss\,\mathbf{\Delta}^1=\Liss_0\mathbf{\Delta}^1$, d'où aussitôt $\Liss\,S=\Liss_0S$, pour tout $S$ dans Cat, ça se saurait! Le contrexemple le plus trivial est de prendre $X=\{0\}\hookrightarrow S=\mathbf{\Delta}^1$, l'inclusion de l'ouvert ‘générique’, immersion ouverte, donc lisse, la fibre $X_1$ est vide! % [sic] s'ensuiverait -> s'ensuivrait, nonstandard verb form

% marginal mark: handwritten sign in the left margin beside the following statement
\textbf{Corollaire 1 rectifié.} \textit{Si $X$ est dans} $\Cat/S$\textit{, pas nécessairement lisse sur $S$, alors le foncteur}
\[\mult_{X/S} : Y\longmapsto Y\times_S X : \Cat/S\longrightarrow \Cat/S\]
\textit{induit un foncteur}
\[\Liss\,S\longrightarrow \Cat/S\]
\textit{compatible avec les localiseurs $W_S^{\mathrm{d}}$, $W_S^{\mathrm{d}}$, donc induit un foncteur}
\[(\Liss\,S)(W_S^{\mathrm{d}})^{-1}\longrightarrow(\Cat/S)(W_S^{\mathrm{d}})^{-1},\]

\msp{155}
\textit{ou encore (compte tenu de l'équivalence induite par l'inclusion}
\[(\underline{\mathrm{Liss}}\,S)(W_S^{\mathrm{d}})^{-1}\xrightarrow{\ \sim\ }(\Cat/S)(W_S^{\mathrm{d}})^{-1}\quad\textit{),}\] % doubt resolved in DER_VI-c_corrections.csv; input note: underline of Liss -- double underline or single underline
\textit{un foncteur, encore noté} $\multhot_{X/S}$,
\[\multhot_{X/S} : \HOT_W(S)\longrightarrow\HOT_W(S).\]

Ce foncteur satisfait donc à la relation
\[\multhot_{X/S}(\hot_{S,W}(Y))\qquad\underbrace{\ \simeq\ }_{\text{isomorphisme canonique}}\qquad\hot_{S,W}(X\times_S Y)\]
% foot: Version 9 novembre 2001
\printedpage{95}
% head: A. Grothendieck : Dérivateurs, Ch. VI, Édition des Universités de Montpellier II et Paris VII
% marginal mark: handwritten sign in the left margin beside the following two lines
pour $Y$ \textit{lisse sur $S$,} et pour $Y$ quelconque dans $\Cat/S$,
\[\multhot_{X/S}(\hot_{S,W}(Y))\qquad\underbrace{\ \simeq\ }_{\text{isomorphisme canonique}}\qquad\hot_{S,W}(X\times_S \Phi_S(Y)).\]

Quand $X$ est lui-même lisse sur $S$, ce foncteur n'est autre que la multiplication par $\hot_{S,W}(X)$ :
\[\multhot_{X/S}=\mult_{\hot_{S,W}(X)}\qquad\text{si $X$ \textit{lisse sur $S$.}}\]

% marginal mark: handwritten sign in the left margin beside the following paragraph
Mais sa valeur pour $\hot_{S,W}(Y)$, quand $Y$ n'est \textit{pas} lisse sur $S$, ne doit pas être confondue avec $\hot_{S,W}(X\times_S Y)$.

Ainsi, si $X$ et $Y$ sont dans $\Cat/S$, considérons
\[X\times_S Y\qquad\text{et}\qquad\boxed{\hot_{S,W}(X\times_S Y).}\] % the frame is printed with rounded corners

\begin{enumerate}
\item[a)] Si ni $X$ ni $Y$ ne sont supposés lisses sur $S$, cet élément ‘n'a aucune interprétation naturelle en termes d'opérations sur les catégories $\HOT_W(S)$’ (\footnote[74]{Sauf dans le cas où l'un des deux soit dans $\underline{\mathrm{Fibhot}}\,S$, cf. annotation page suivante.}).

\msp{156}

\item[b)] Si $X$ et $Y$ sont lisses sur $S$, cet élément s'interprète comme
\[\hot_{S,W}(X)\times_S\hot_{S,W}(Y)\]
(\footnote[75]{Mais aussi si l'\textit{un} des deux éléments est dans $\underline{\mathrm{Fibhot}}\,S$ (et moyennant l'axiome L6 bis sur $W$).}).

\item[c)] Si l'un des deux $X$, $Y$ est lisse sur $S$, sans rien supposer sur l'autre, on a une interprétation en termes de $\mult_X$ (si $Y$ lisse), ou $\mult_Y$ (si $X$ lisse), notamment
\[\begin{aligned}
\hot_{S,W}(X\times_S Y)\ &\simeq\ \multhot_X(\hot_{S,W}(Y))\quad\text{si $Y$ lisse sur $S$}\\
&\simeq\ \multhot_Y(\hot_{S,W}(X))\quad\text{si $X$ lisse sur $S$.}
\end{aligned}\]
\end{enumerate}

\textbf{Corollaire 2.} \textit{Soit $f : S'\longrightarrow S$ une flèche dans} Cat\textit{, d'où $f_!$, $f^*$ entre} $\HOT_W(S')$ \textit{et} $\HOT_W(S)$\textit{. On a la relation}
\[f_!f^*(\xi)\simeq\multhot_{S'}(\xi),\]
\textit{donc si $S'$ est lisse sur $S$,}
\[f_!f^*(\xi)\simeq\hot_{S,W}(S')\times_S\xi.\]

Nous sommes maintenant en mesure de vérifier les axiomes Der 4$'$b, Der 5$'$ relatifs aux images directes (I, p.~67, 68), et de plus Der 3b (I, p.~111). L'axiome Der 4$'$ n'est autre que l'existence des adjoints à gauche $f_!$ des $f^*$. L'axiome Der 5$'$ permet le ‘calcul’ des fibres d'un objet $f_!(\xi)$. Il dit qu'une certaine
% foot: Version 9 novembre 2001
\printedpage{96}
% head: A. Grothendieck : Dérivateurs, Ch. VI, Édition des Universités de Montpellier II et Paris VII
\msp{157}
flèche du changement de base, relatif au diagramme commutatif dans Cat
\[\begin{tikzcd}[row sep=large]
S' \arrow[d,"f"'] & s\backslash S' \arrow[l,"q"'] \arrow[d,"f'"] \\
S & e\simeq\{s\} \arrow[l,"p"']
\end{tikzcd}\qquad(s\in\mathrm{Ob}\,S),\]
savoir
\[p^*(f_!(\xi'))\longrightarrow f'_!(q^*(\xi'))\]
est un isomorphisme (pour $\xi'\in\mathbf{D}(S')$). Soit donc
\[\xi'\in\mathrm{Ob}\,\HOT_W(S'),\qquad\xi'=\hot_{S',W}(X')\]
avec $X'$ dans $\Cat/S'$. Comme
\[f_!(\xi')=\hot_{S,W}(X'),\]
et que $X'$ n'a aucune raison d'être lisse sur $S$ (si on ne suppose pas $S'$ lui-même lisse sur $S$, auquel cas il pourrait être utile de prendre $X'$ lisse sur $S'$, donc aussi sur $S$), pour calculer $p^*(f_!(\xi'))$ il n'est pas question de prendre la fibre $X'_s$, mais il faut prendre $s\backslash X'=\Phi_S(X')_s$ :
\[p^*(f_!(\xi'))\simeq\hot_W(s\backslash X').\]
\[\begin{tikzcd}[row sep=large]
X' \arrow[d] & s\backslash X' \arrow[l] \arrow[d] \\
S' \arrow[d] & s\backslash S' \arrow[l,"q"'] \arrow[d] \\
S & s\backslash S \arrow[l]
\end{tikzcd}\]

Considérons le diagramme de changement de base par $s\backslash S\longrightarrow S$ ci-contre \texttt{[i.e. ci-dessus]}, à carrés cartésiens. On voit d'abord que $q^*(\xi')$ n'est autre que la classe dans $\HOT_W(s\backslash S')$ de l'image inverse de $X'$ par le changement de base $s\backslash S'\longrightarrow S'$. Pour l'assurer, il suffirait de prendre $X'$ lisse sur $S'$ (ce qui
\msp{158}
est licite) - alors la description de l'image inverse est OK pour tout changement de base. Mais on a vu que dans le cas d'un co-isomorphisme local comme $s\backslash S'\longrightarrow S'$, il n'y a pas même besoin de cet artifice - la formule est OK sans restriction de lissité. Ainsi
\[q^*(\xi')=\hot_{s\backslash S',W}(s\backslash X'),\]
% foot: Version 9 novembre 2001
\printedpage{97}
% head: A. Grothendieck : Dérivateurs, Ch. VI, Édition des Universités de Montpellier II et Paris VII
d'où
\[f'_!(q^*(\xi'))\simeq\hot_{S,W}(s\backslash X')=\hot_W(s\backslash X').\]
Donc on a bien trouvé
\[p^*f_!\simeq f'_!q^*.\]

On peut maintenant prouver aussi Der 3b. On a vu qu'une fois qu'on dispose des adjoints à gauche $f_!$ des $f^*$, cet axiome s'énonce sous \textit{forme homologique} ainsi : Si
\[f : S'\longrightarrow S\]
est dans $W$, il faut prouver que pour tout
\[\xi\in\HOT_W(e)\overset{\text{déf}}{=}\mathrm{Hot}_W,\]
la flèche correspondante
\[(*)\qquad p'_!p'^*(\xi)\longrightarrow p_!p^*(\xi)\]
est \texttt{[un]} isomorphisme dans $\mathrm{Hot}_W$, où $p : S\longrightarrow e$, $p' : S'\longrightarrow e$ sont les flèches structurales. Or comme tout objet de Cat est lisse sur $e$, le corollaire 2 précédent permet d'identifier cette flèche $(*)$ à
\[\eta'\times_S\xi\longrightarrow\eta\times_S\xi\]
\msp{159}
déduite de
\[\hot_W(f) :\quad\underbrace{\ \eta'\ }_{\overset{\text{déf}}{=}\,\hot_W(S')}\quad\longrightarrow\quad\underbrace{\ \eta\ }_{\overset{\text{déf}}{=}\,\hot_W(S)}\quad.\]
Comme celle-ci est \texttt{[un]} isomorphisme par hypothèse, on gagne.

Donc on a prouvé

\textbf{Théorème 3.} \textit{Le prédérivateur} HOT \textit{satisfait aux axiomes} Der 1, Der 2 (\footnote[76]{\textit{pas prouvé}, le ‘résultat’ auquel je pensais est canulé, cf. complément proposition 5, page 164.}) \textit{(pour mémoire), et} Der 3b (\footnote[77]{Et même la variante la plus forte Der 3b$'_2$ de Der 3, cf. démonstration page suivante. \textbf{NB} Pour Der 3b$'_2$, on a utilisé l'axiome L6 sur $W$ - et c'est essentiel. On voit que la validité de Der 3b$'_2$ implique L6, et même la forme plus forte L6 bis, et même si on suppose seulement Der 3b$'_1$, i.e. $f_0^*$ pleinement fidèle.}), Der 4$'$, Der 5$'$\textit{.} % [sic] canulé (footnote 76) -> annulé (?), not a French word

On a envie de prouver aussi les variantes renforcées de Der 3, la forme la plus forte étant celle-ci (Der 3b$'_2$) : si $f\in W$, alors $f^*$ induit une \textit{équivalence} de catégories
\[f_0^* : \HOT^{\mathrm{lc}}_W(S)\longrightarrow\HOT^{\mathrm{lc}}_W(S').\]

La pleine fidélité du foncteur $f_0^*$ doit résulter de la variante naturelle Der 3b$''_{\mathrm{coh}}$ de Der 3b, une fois qu'on aura développé le formalisme des Hom internes dans les $\HOT_W(S)$. Mais pour donner un sens à Der 3b$''_{\mathrm{coh}}$, on a besoin des foncteurs \textit{cohomologie,} i.e. de $p_*$ et $p'_*$ (au lieu de $p_!$, $p'_!$) pour HOT, donc c'est encore prématuré pour l'instant. La construction des $f_*$ et celle des Hom internes se fera d'ailleurs sûrement $\pm$ conjointement.
% foot: Version 9 novembre 2001

\printedpage{98}
% head: A. Grothendieck : Dérivateurs, Ch. VI, Édition des Universités de Montpellier II et Paris VII
\msp{160}
Mais je m'aperçois à l'instant qu'on trouve quasiment ‘à l'œil’ le fait que $f_0^*$ soit bien une \textit{équivalence} de catégories, et que $f_!^0$ soit l'équivalence quasi-inverse, en utilisant l'adjonction des foncteurs $f_!^0$, $f_0^*$, et se rappelant que pour prouver que ce sont des équivalences quasi-inverses l'une de l'autre, il revient au même de montrer que les morphismes d'adjonction
\[\begin{aligned}
\mathrm{id}&\quad\xrightarrow{\ \alpha_0\ }\quad f_0^*f_!^0\\
\mathrm{id}&\quad\xleftarrow{\ \beta_0\ }\quad f_!^0f_0^*
\end{aligned}\]
sont \texttt{[des]} isomorphismes. Je vais essayer de décrire ces flèches d'adjonction dans le cas général (sans supposer $f\in W$), puis de montrer qu'elles sont dans $W$ si $f$ l'est.

\begin{enumerate}
\item[\texttt{[1)}] \texttt{Soit]} $\xi=\hot^{\mathrm{lc}}_{S,W}(X)$ dans $\HOT_W(S)$. On peut supposer $X\in\mathrm{Ob}\,\underline{\mathrm{Parf}}\,S$. Comme \texttt{[‘Comme’ barré]} $X$ lisse sur $S$, on a
\[f_0^*(\xi)=\hot^{\mathrm{lc}}_{S',W}(X\times_S S'),\qquad f_!^0f_0^*(\xi)=\hot^{\mathrm{lc}}_{S,W}(X\times_S S').\]
Donc
\[\beta_{0,\xi} : f_!^0f_0^*(\xi)\longrightarrow\xi\]
est une flèche
\[X\times_S S'\longrightarrow X\qquad\text{dans }(\Cat/S)(W_S)^{-1}\]
et je présume que c'est $\mathrm{pr}_1$ (\footnote[78]{\textbf{NB} On voit que l'axiome L6 est $\pm$ \textit{nécessaire} pour que le résultat sur $f_0^*$, $f_!^0$ soit valable.}). Est-ce dans $W$ si $f : S'\longrightarrow S$ l'est? Oui (car $X$ est parfait

\msp{161}

sur $S$), sous réserve que $W$ satisfait l'axiome L6 (cf. page 43).

\item[2)] Soit
\[\xi'=\hot^{\mathrm{lc}}_{S',W}(X')\qquad\text{dans }\HOT_W(S'),\]
on a
\[f_!^0(\xi')=\hot^{\mathrm{lc}}_{S,W}(X'),\]
mais pour expliciter $f_0^*$ de cet élément, comme $X'$ n'a aucune raison d'être un Hot-fibré sur $S$ (auquel cas il suffirait de prendre $f^*(X')\overset{\text{déf}}{=}X\times_S S'$), il faut d'abord prendre $C_S(X')$, puis appliquer $f^*$ :
\[f_0^*f_!^0(X')=\hot^{\mathrm{lc}}_{S',W}(C_S(X')\times_S S').\]

Il faut donc définir
\[X'\xrightarrow{\beta_{0,X'}}C_S(X')\times_S S'\qquad\text{dans }(\Cat/S')W_{S'}^{-1}.\]

Or on trouve une flèche canonique dans $\Cat/S'$, déduite du $S$-morphisme $i_{X'}$ : $X'\longrightarrow C_S(X')$, i.e. du carré commutatif en traits pleins
\end{enumerate}
% foot: Version 9 novembre 2001

\end{document}
